% Part: computability
% Chapter: computability-theory
% Section: non-comp-set

\documentclass[../../../include/open-logic-section]{subfiles}

\begin{document}

\olfileid{cmp}{thy}{ncp}
\olsection{অগণনসাধ্য সেট আছে}


উপরে আমরা দেখেছি, প্রত্যেক গণনসাধ্য সেট গণনসাধ্যভাবে তালিকায়নযোগ্য।
বিপরীত বক্তব্যটি কি সত্য? নিচের ফল দেখায়, সাধারণভাবে তা সত্য নয়।

\begin{thm}\ollabel{thm:K-0}
ধরা যাক, $K_0$ হলো সেট $\Setabs{\tuple{e, x}}{\cfind{e}(x) \fdefined}$।
তাহলে $K_0$ গণনসাধ্যভাবে তালিকায়নযোগ্য, কিন্তু গণনসাধ্য নয়।
\end{thm}

\begin{proof}
$K_0$ যে গণনসাধ্যভাবে তালিকায়নযোগ্য তা দেখতে লক্ষ করুন, এটি নিচের
সংজ্ঞায় নির্ধারিত অপেক্ষক~$f$-এর সংজ্ঞাক্ষেত্র:
\[
f(z) = \umin{y}{(\len{z} = 2 \land T((z)_0, (z)_1, y))}.
\]
কারণ $\cfind{e}(x)$ সংজ্ঞায়িত হলে $f(\tuple{e, x})$ একটি থামা
গণনা-অনুক্রম খুঁজে পায়; $\cfind{e}(x)$ অসংজ্ঞায়িত হলে
$f(\tuple{e, x})$-ও অসংজ্ঞায়িত; আর $z$ যদি কোনো জোড়াকেই সংকেতায়িত না
করে, তবে $f(z)$-ও অসংজ্ঞায়িত।

$K_0$ গণনসাধ্য নয়—এই তথ্যটি শুধু থামার সমস্যার অনির্ণেয়তা,
\olref[hlt]{thm:halting-problem}।
\end{proof}

সেট $K_0$ হলো এমন জোড়া $\tuple{e,x}$-এর সেট, যাদের জন্য
$\cfind{e}(x) \fdefined$; অর্থাৎ $\tuple{e,x} \in K_0$ যদি এবং কেবল যদি
$\cfind{e}$ নিবেশ~$x$-এ সংজ্ঞায়িত (থামে)। তাই একে “থামা-সেট”-ও বলা হয়।
সেট $K = \Setabs{e}{\cfind{e}(e) \fdefined}$ হলো “স্ব-থামা সেট”। একে
প্রায়ই একটি আদর্শ অনির্ণেয় সেট হিসেবে ব্যবহার করা হয়।

\begin{thm}\ollabel{thm:K}
স্ব-থামা সেট $K = \Setabs{e}{\cfind{e}(e) \fdefined}$ !!{c.e.}, কিন্তু
নির্ণেয় নয়।
\end{thm}

\begin{proof}
  ধরা যাক, $K$ নির্ণেয়, অর্থাৎ তার চরিত্রসূচক অপেক্ষক $\Char{K}$
  গণনসাধ্য। ধরি,
  \[d(e) = \begin{cases}
  1 & \text{যদি\/ $\Char{K}(e) = 0$}\\
  \fundefined & \text{অন্যথায়।}
  \end{cases}
  \]
  ধরা যাক, $k$ হলো~$d$-এর সূচক, অর্থাৎ $d \simeq \cfind{k}$। তাহলে
  $d(k) \simeq \cfind{k}(k)$। কিন্তু~$d$-এর সংজ্ঞা থেকে
  $d(k) \fdefined$ যদি এবং কেবল যদি $\cfind{k}(k) \fundefined$; তাই
  বিরোধ পাওয়া যায়।

  $K$ হলো $f(x) = \umin{y}{T(x,x,y)}$-এর সংজ্ঞাক্ষেত্র; অতএব সেটি
  !!{c.e.}
\end{proof}

\end{document}
